% Generated from validated Special post-draft synthesis. Do not hand-edit.
\section*{Monthly Signals}
\addcontentsline{toc}{section}{Monthly Signals}
\smallskip
\noindent \textbf{Codingから継続実行へ} 複数agent、長時間task、tool useが同時に前景化し、coding能力だけでなく仕事を持ち続ける実行系そのものが競争軸になった。 \hfill{\footnotesize p.~\pageref{pkg:agentic-engineering-frontier} / p.~\pageref{pkg:open-model-agent-stack}}
\par
\smallskip
\noindent \textbf{Interactionの時間とmodalが広がる} streaming speechとimage generationが、単発のtext応答とは異なる連続的なinteractionを実装層へ押し広げた。 \hfill{\footnotesize p.~\pageref{pkg:realtime-multimodal-interaction}}
\par
\smallskip
\noindent \textbf{科学推論は独立したfrontier} science、research、engineering向けのspecialized reasoningは、汎用agent競争とは別に、workloadへ推論を特化する設計軸を示した。 \hfill{\footnotesize p.~\pageref{pkg:specialized-reasoning-science}}
\par
\smallskip
\noindent \textbf{Servingが能力を成立させる} realtime audio、tool protocol、新しいmodel familyへの対応がruntimeへ集中し、servingは発表後の最適化ではなくmodel capabilityを成立させる機能層になった。 \hfill{\footnotesize p.~\pageref{pkg:serving-stack-catches-up}}
\par
\smallskip
\noindent \textbf{後続月を読むための基準点} 2月に立ち上がったagent実行、runtime、安全な制御、multimodal interaction、specialized reasoningという観察軸を、後続月の変化を比較する基準として残す。 \hfill{\footnotesize p.~\pageref{pkg:closing-synthesis}}
\par
\medskip
\begin{claimboundary}[Retrospective scope]
本号は2026年2月を後日確認可能になった一次情報も用いて再構成するRetrospective Specialである。Coverage Periodと制作時点を同一視せず、vendor / project / author claimの境界は各記事で明示する。
\end{claimboundary}
\medskip
\tableofcontents
